\section{List of acronyms} % non-titlecase
%\section{List of Acronyms} % titlecase
% titlecase should be either consistently used or consistently not be used.
% mixing titlecase with non-titlecase should be avoided.
\label{sec:acronym}

\begin{acronym}[PCMCIA] % Put the longest acronym in this list into the []
	\acro{DuT}{Device under Test\acroextra{, a device that is investigated during experiments}}%
	\acro{XDP}{EXpress Data Path\acroextra{, a programmable hook within the Linux kernel allowing for early packet processing before the network stack.}}
	\acro{PTP}{Precision Time Protocol\acroextra{, an IEEE standard that enables nanosecond-level clock synchronization across computer networks.}}
	\acro{HTTP/3}{Hypertext Transfer Protocol 3\acroextra{, the latest version of the Hypertext Transfer Protocol, which leverages QUIC over UDP to eliminate head-of-line blocking.}}
	\acro{BPF}{Berkeley Packet Filter\acroextra{, in-kernel virtual machine that allows sandboxed bytecode to be executed at various hook points.}}
	\acro{eBPF}{extended Berkeley Packet Filter\acroextra{, an extended, in-kernel bytecode language and virtual machine, which expands on cBPF's capabilities to allow for programmable, event-driven execution across the kernel's subsystems.}}
	\acro{cBPF}{classic Berkeley Packet Filter\acroextra{, a minimalistic, in-kernel bytecode language and virtual machine primarily used for network packet inspection and filtering.}}
	\acro{DDoS}{Distributed Denial of Service\acroextra{, a malicious attempt to disrupt the normal traffic of a targeted service by flooding it with internet traffic from multiple sources.}}
	\acro{RISC}{Reduced Instruction Set Computer\acroextra{, an instruction set architecture that prioritizes simpler, fixed-length instructions for faster execution and simpler processor design.}}
	\acro{JIT}{Just in Time\acroextra{, a method where code is translated from an intermediate representation into native machine code during runtime, immediately before execution, to improve performance.}}
	\acro{NIC}{Network Interface Card\acroextra{, a hardware component that connects a computer to a computer network, allowing it to send and receive data.}}
	\acro{RSS}{Receive Side Scaling\acroextra{, a network driver feature that distributes the processing of incoming network packets across multiple CPU cores to improve network throughput and reduce CPU bottlenecks.}}
	\acro{XSK}{XDP Socket\acroextra{, a Linux kernel socket and part of the AF\_XDP address family, which allows user-space applications to directly access raw network packets.}}
	\acro{DMA}{Direct Memory Access\acroextra{, a feature that allows hardware components to directly access system memory for reading and writing data without requiring the CPU's intervention.}}
\end{acronym}
